\documentclass[11pt]{article}
\input{packages.tex}
\usepackage{ulem}
\DeclareMathOperator{\sign}{sign}
\DeclareMathOperator{\Imm}{Im}
\DeclareMathOperator{\DIF}{DIF}
\DeclareMathOperator{\rg}{rg}
\DeclareMathOperator{\dist}{dist}

% Настройка колонтитулов
\pagestyle{fancy}
\fancyhf{} % Очистить все поля
\fancyhead[L]{Злых людей нет на свете, есть только люди несчастливые}
\fancyhead[R]{\href{https://t.me/ceagest}{Автор}}
\fancyfoot[C]{\thepage}               % Справа внизу

\setlength{\headheight}{13.6pt}

\hypersetup{
    colorlinks=true,
    linkcolor=blue,
    citecolor=blue,
    filecolor=blue,
    urlcolor=blue,
}

\begin{document}

\begin{titlepage}
    \centering
    \vspace*{\fill}  % заполняет пространство сверху

    {\Huge Лекция 15. Следствия из критерия Лебега. Множество Кантора. Критерий интегрируемости в терминах интегральных сумм Римана. Линейность интеграла Римана.
    \par}

    \vspace*{\fill}  % заполняет пространство снизу
\end{titlepage}

\setcounter{page}{1}

\section{Следствия из критерия Лебега}

\begin{theorem}
    Если $f \in C([a, b])$, то $f$ интегрируема по Риману на $[a, b]$.
\end{theorem}

\begin{proof}
    Действительно, $f$ ограниченна на $[a, b]$ по теореме Вейерштрасса, так как она непрерывна на этом отрезке. Кроме того, всё также в силу непрерывности множество точек разрыва функции $f$ на $[a, b]$ пусто, а значит оно имеет лебегову меру ноль. Тогда по критерию Лебега немедленно получаем требуемое.
\end{proof}

\begin{theorem}
    Если $f$ --- монотонна на $[a, b]$, то она интегрируема по Риману на отрезке $[a, b]$.
\end{theorem}

\begin{proof}
    Ясен пень, $f$ ограниченна на $[a, b]$, так как 
    \[f(x) \in [\min\{f(a), f(b)\}, \max\{f(a), f(b)\}] \quad \forall x \in [a, b]\]
    При этом у монотонной функции по доказанному в прошлом семестре могут быть лишь точки разрыва первого рода, а множество этих точек разрыва не более, чем счётно. Значит множество точек разрыва $f$ на $[a, b]$ не более, чем счётно, откуда сразу получаем, что это множество имеет лебегову меру ноль. Теперь с радостью применяем критерий Лебега и получаем удовольствие.
\end{proof}

\begin{theorem}
    Если $f$ интегрируема по Риману на $[a, b]$, то она интегрируема по Риману на любом отрезке $[a', b'] \subset [a, b]$.
\end{theorem}

\begin{proof}
    Так как $f$ интегрируема по Риману на $[a, b]$, то по критерию Лебега она ограниченна на $[a, b]$ и её множество точек разрыва на этом отрезке имеет лебегову меру ноль. Тогда понятно, что $f$ также ограниченна и на любом отрезке $[a', b'] \subset [a, b]$, и, кроме того, множество точек разрыва $f$ на $[a', b']$ очевидным образом является подмножеством множества её точек разрыва на отрезке $[a, b]$, а значит это множество также имеет лебегову меру ноль. Осталось лишь воспользоваться критерием Лебега для $[a', b']$ и получить требуемое.
\end{proof}

\begin{definition}
    Пусть $f \colon [a, b] \to \mathbb{R}$. Тогда $\mathcal{D}(f, [a, b])$ --- множество точек разрыва функции $f$ на отрезке $[a, b]$, а $\mathcal{C}(f, [a, b])$ --- множество точек непрерывности функции $f$ на отрезке $[a, b]$.
\end{definition}

\begin{theorem}
    Если $f, g \in \mathcal{R}([a, b])$, то $f g \in \mathcal{R}([a, b])$.
\end{theorem}

\begin{proof}
    Если $f, g \in \mathcal{R}([a, b])$, то по критерию Лебега $f$ и $g$ ограниченны на $[a, b]$. Тогда ясно, что и $f g$ ограниченна на $[a, b]$. Далее заметим, что 
    \begin{equation*} \label{eq:th:1.4} \tag{$\ast$}
        \mathcal{C}(f, [a, b]) \cap \mathcal{C}(g, [a, b]) \subset \mathcal{C}(fg, [a, b])
    \end{equation*}
    В самом деле, из непрерывности функций $f$ и $g$ в произвольной точке отрезка $[a, b]$ по теореме о произведении непрерывных функций следует непрерывность в этой точке функции $fg$. Это и означает, что множество $\mathcal{C}(fg, [a, b])$ содержит уж по крайней мере те точки отрезка $[a, b]$, в которых одновременно непрерывны обе рассматриваемые функции. Тогда, взяв дополнение до отрезка $[a, b]$ от обеих частей включения (\ref{eq:th:1.4}) и вспомнив, что дополнение к пересечению есть объединение дополнений, получаем:
    \[\mathcal{D}(fg, [a, b]) \subset \mathcal{D}(f, [a, b]) \cup \mathcal{D}(g, [a, b])\]
    Но по критерию Лебега множества $\mathcal{D}(f, [a, b])$ и $\mathcal{D}(g, [a, b])$ имеют лебегову меру ноль, а значит и их объединение имеет лебегову меру ноль по доказанному в предыдущей лекции (на самом деле был показан более общий факт о объединении счётного числа множеств лебеговой меры ноль). При этом очевидно, что подмножество $\mathcal{D}(fg, [a, b])$ множества лебеговой меры ноль само имеет лебегову меру ноль. Тогда, собирая два наших вывода в одно целое прекрасное утверждение про функцию $fg$ на отрезке $[a, b]$, по критерию Лебега получаем то, что и требовалось.
\end{proof}

\begin{definition}
    Функция $f \colon [a, b] \to \mathbb{R}$ называется \textbf{кусочно-непрерывной} на отрезке $[a, b]$, если существует разбиение 
    \[T = \{x_i\}_{i = 0}^{N_T}, \quad a = x_0 < \ldots < x_{N_T} = b, \]
    такое, что:
    \begin{enumerate}
        \item $f \in C((x_{i - 1}, x_i)) \quad \forall i \in \{1, \ldots, N_T\}$
        \item в точках $\{x_i\}$ существуют односторонние пределы функции $f$
    \end{enumerate}
\end{definition}

\begin{figure}[h]
\centering
\includegraphics[width=0.7\linewidth]{pictures/Кусочно непрерывная функция.jpg}
\end{figure}

\begin{theorem}
    Если $f$ кусочно-непрерывна на $[a, b]$, то $f \in \mathcal{R}([a, b])$.
\end{theorem}

\begin{proof}
    Так как $f$ кусочно-непрерывна на $[a, b]$, то по определению множество $\mathcal{D}(f, [a, b])$ конечно. Тогда это множество, очевидно, имеет лебегову меру ноль. Остаётся доказать ограниченность $f$ на $[a, b]$. Для этого достаточно заметить, что $f$ ограниченна на любом интервале вида $(x_{i - 1}, x_i)$. В самом деле, рассмотрим функцию $\widetilde{f}$ для каждого интервала $(x_{i - 1}, x_i)$ такую, что:
    \[\widetilde{f}(x_{i - 1}) := f(x_{i - 1} + 0)\]
    \[\widetilde{f}(x_{i}) := f(x_{i} - 0)\]
    \[\widetilde{f}(x) := f(x) \quad \forall x \in (x_{i - 1}, x_i)\]
    Несложно видеть, что $\widetilde{f} \in C([x_{i - 1}, x_i])$, откуда по теореме Вейерштрасса функция $\widetilde{f}$ ограниченна на отрезке $[x_{i - 1}, x_i]$. \\
    При этом функция $\widetilde{f}$ по определению совпадает с функцией $f$ на $(x_{i - 1}, x_i)$ и, очевидным образом, является ограниченной также на интервале $(x_{i - 1}, x_i)$, а значит уже $f$ является ограниченной на $(x_{i - 1}, x_i)$, что демонстрирует справедливость прекрасного замечания выше. \\
    Осталось лишь вспомнить, что количество интервалов соответствующего вида и точек разрыва конечно по определению кусочно-непрерывной функции, а значит мы можем положить границами $f$ на всём отрезке $[a, b]$ максимум и минимум среди соответственно конечных супремумов и инфимумов на интервалах разбиения и значениях функции в точках разрыва. Это и означает ограниченность функции $f$ на отрезке $[a, b]$. \\ 
    Последним действием в этом доказательстве будет радостное объединение двух уже знакомых до боли доказанных фактов про функцию $f$ на отрезке $[a, b]$ и упоминание чудесного заклинания, называемого критерием Лебега.
\end{proof}

\begin{lemma}
    Пусть $f, g \colon [a, b] \to \mathbb{R}$. Пусть $f$ и $g$ отличаются лишь в конечном числе точек. Тогда либо $f$ и $g$ одновременно интегрируемы по Риману на $[a, b]$, либо $f$ и $g$ одновременно не интегрируемы по Риману на $[a, b]$.
\end{lemma}

\begin{proof}
    В самом деле, если одна из функций ограниченна на $[a, b]$ и отличается от другой лишь в конечном числе точек, то и другая, очевидным образом, будет являться ограниченной на $[a, b]$. Если для одной из функций множество её точек разрыва на $[a, b]$ имеет лебегову меру ноль и эта функция отличается от другой на конечное число точек, то и для другой функции множество точек разрыва на $[a, b]$ будет иметь лебегову меру ноль (более подробное доказательство этого чудесного утверждения есть в начале следующей лекции). По совокупности этих фактов из критерия Лебега незамедлительно получаем требуемое.
\end{proof}

\begin{question}
    Пусть $f, g \colon [a, b] \to \mathbb{R}$ ограниченны на $[a, b]$. Пусть $f$ и $g$ отличаются в точках счётного множества. Верно ли, что $f$ интегрируема на $[a, b]$ тогда и только тогда, когда $g$ интегрируема на $[a, b]$?
\end{question}

\begin{answer}
    Нет, неверно, так как изменение функции на счётном множестве может привести к добавлению несчётного множества точек разрыва.
\end{answer}

\begin{example}
    Пусть $f \equiv 0$ на отрезке $[0, 1]$. Пусть
    \begin{equation*}
        g(x) = 
        \begin{cases}
            1, & x \in \mathbb{Q} \cap [0, 1] \\
            0, & x \in (\mathbb{R} \setminus \mathbb{Q}) \cap [0, 1]
        \end{cases}
    \end{equation*}
    Очевидно, что 
    \[\mathcal{D}(f, [0, 1]) = \varnothing \quad \text{ и } \quad \mathcal{D}(g, [0, 1]) = [0, 1]\]
    Увидим, что множество $[0, 1]$ не имеет лебегову меру ноль. Действительно, если бы от противного это было так, то по определению существовал бы счётный набор интервалов, покрывающих этот отрезок. Но интервал является открытым множеством, а значит по лемме Гейне-Бореля из этого покрытия можно выделить конечное подпокрытие. Совершенно очевидно, что сумма длин интервалов этого конечного покрытия подпирается снизу длиной отрезка $[0, 1]$, откуда и сумма длин всех интервалов исходного покрытия уж точно не меньше единицы. Но эта сумма длин должна быть сколь угодно мала по определению множества лебеговой меры ноль, откуда получаем противоречие. \\
    Итого по критерию Лебега получаем, что $f \in \mathcal{R}([0, 1])$ и $g \not \in \mathcal{R}([0, 1])$, а значит рассматриваемое утверждение не справедливо.
\end{example}

\begin{remark}
    С другой стороны, есть примеры, когда две функции отличаются на несчётном множестве, но при этом обе являются интегрируемыми по Риману.
\end{remark}

\begin{example}[Множество Кантора]
    Рассмотрим прекрасное множество, называемое \textbf{множеством Кантора} (или, как говорят более конкретно в данном случае англичане, \textbf{middle-third Cantor set}, так как на самом деле есть целое семейство множеств Кантора, а мы рассматриваем одно конкретное). Его построение происходит итеративно (на очередной итерации над каждым из отрезков в объединении проводится следующая процедура: данный отрезок разбивается на три равные части, серединка удаляется, а две оставшееся части объединяются):
    \[K_0 = [0, 1], \quad K_1 = \left[0, \frac{1}{3}\right] \cup \left[\frac{2}{3}, 1\right], \quad K_2 = \left[0, \frac{1}{9}\right] \cup \left[\frac{2}{9}, \frac{1}{3}\right] \cup \left[\frac{2}{3}, \frac{7}{9}\right] \cup \left[\frac{8}{9}, 1\right], \ldots\]
    \begin{figure}[h]
    \centering
    \includegraphics[width=0.6\linewidth]{pictures/Множество Кантора.jpg}
    \end{figure} \leavevmode
    \\
    Несложно видеть, что на каждой итерации проводится процедура самоподобия, поэтому множества Кантора является простейшим фракталом. Также очевидно, что на $n$-ом шаге $K_n$ состоит из $2^n$ конгруэнтных отрезков длины $3^{-n}$. Итоговое же множество есть
    \[K := \bigcap_{n = 0}^{\infty} K_n\]
    Поймём теперь, как же устроено это загадочное множество $K$. Легко видеть, что $\forall n \in \mathbb{N}$ множество $K_n$ --- замкнуто, а значит замкнуто и множество $K$ как счётное пересечение замкнутых. Кроме того, ясно, что $K$ точно содержит рациональные точки, представимые дробями
    \begin{equation*} \label{eq:ex:1.2} \tag{$\ast$}
        \left\{\frac{p}{3^n}\right\}, \; \; n \in \mathbb{N}, \; \; p \in \{0, \ldots, 3^n\}        
    \end{equation*}
    Более того, $K$ также содержит всевозможные пределы точек вида (\ref{eq:ex:1.2}) в силу замкнутости множества Кантора (под всевозможными пределами точек понимаются всевозможные пределы последовательностей, состоящих из точек такого вида). \\
    Теперь увидим, что множество $K$ имеет лебегову меру ноль. Каждый отрезок из $K_n$ погрузим в интервал длины $\frac{3}{2} \left(\frac{1}{3}\right)^n$. Тогда всё $K_n$ окажется покрыто этими интервалами. Суммарная же длина данных интервалов есть $\frac{3}{2} \left(\frac{2}{3}\right)^n \to 0, \; n \to \infty$. Но каждое покрытие $K_n$ является покрытием и для канторова множества, так как $\forall m > n$ множество $K_m$ строго содержится в множестве $K_n$. Тогда для любого $\varepsilon > 0$ существует счётный набор интервалов, служащий покрытием для такого множества $K_n$, что $\frac{3}{2} \left(\frac{2}{3}\right)^n < \varepsilon$, который покрывает $K$, откуда по определению множества лебеговой меры ноль получаем, что $K$ имеет лебегову меру ноль. \\
    Докажем далее, что множество $K$ несчётно (а точнее, имеет мощность континуума). Для каждого $n \in \mathbb{N}$ пронумеруем слева направо отрезки множества $K_n$ от $1$ до $2^n$ и обозначим $m$-ый отрезок за $I_{n, m}$. Заметим, что 
    \[\forall x \in K = \bigcap_{n = 0}^{\infty} K_n = \bigcap_{n = 0}^{\infty} \bigcup_{m = 1}^{2^n} I_{n, m}\]
    \[\exists \{I_{n, m(n)}\}_{n = 1}^{\infty} \text{ --- стягивающаяся последовательность вложенных отрезков такая, что }\]
    \[\bigcap_{n = 1}^{\infty} I_{n, m(n)} = x\]
    Действительно, коль скоро $x \in K$, для любого $n \in \mathbb{N}$ он содержится в некотором отрезке $I_{n, m(n)}$ множества $K_n$. Отрезки полученной последовательности являются вложенными просто из способа построения канторова множества, а сама последовательность является стягивающейся, так как длина отрезка $I_{n, m(n)}$ равна $3^{-n} \to 0, \; n \to \infty$. При этом все отрезки по построению содержат точку $x$, а по теореме Кантора о вложенных отрезках все эти товарищи имеют единственную общую точку. Тогда эта точка и есть $x$, что и требовалось. \\
    Таким образом, каждой точке $x \in K$ взаимно-однозначно сопоставлена бесконечная последовательность из нулей и единиц, задающая некоторый путь от корня в бесконечном двоичном дереве, которое строится похожим на процедуру построения канторова множества образом: корень дерева соответствует отрезку $I_{0, 1}$, а каждая вершина имеет двух детей, соответствующих левому и правому отрезку, на которые на очередной итерации разбился родитель, и при этом на ребре в левого ребёнка пишется ноль, а на ребре в правого --- единица. Если в процессе прохождения дерева от корня взяли в последовательность левый отрезок разбиения текущего отрезка, то от текущей вершины проходим по нулевому ребру, а иначе --- по единичному.
    \begin{figure}[h]
    \centering
    \includegraphics[width=0.6\linewidth]{pictures/Бесконесное двоичное дерево.jpg}
    \end{figure} \leavevmode
    \\
    Но множество всех счётных наборов из нулей и единиц имеет мощность континуум, а значит и множество $K$ имеет мощность континуум, что и требовалось. \\
    Осталось лишь, пользуясь доказанными нами фактами, предъявить две функции, отличающиеся на несчётном множестве, но являющиеся интегрируемыми по Риману. Пусть $f \equiv 0$ на $[0, 1]$. Пусть 
    \begin{equation*}
        g(x) = 
        \chi_K (x) =
        \begin{cases}
            1, & x \in K \\
            0, & x \in [0, 1] \setminus K
        \end{cases}
    \end{equation*}
    Так как $K$ несчётно, то $f$ и $g$ отличаются в несчётном множестве точек. При этом множество $\mathcal{D}(f, [0, 1]) = \varnothing$ имеет лебегову меру ноль и множество $\mathcal{D}(g, [0, 1])$ имеет лебегову меру ноль, так как является подмножеством канторова множества, имеющего по доказанному лебегову меру ноль (точки разрыва функции $g$ могут лежать лишь в канторовом множестве, так как дополнение до множества Кантора есть счётное объединение интервалов, а значит это дополнение является открытым множеством, в которое точки входят по самому его определению с некоторой окрестностью, и при этом на всей этой окрестности $g$ тождественно равна нулю, откуда получаем её непрерывность в каждой точке дополнения до $K$). Кроме того, тривиально видеть, что $f$ и $g$ ограниченны на $[0, 1]$. Тогда по критерию Лебега немедленно получаем, что $f, g \in \mathcal{R}([0, 1])$.
\end{example}

\section{Критерий интегрируемости в терминах интегральных сумм \\ Римана}

\begin{definition}
    Пусть $T = \{x_i\}_{i = 0}^{N_T}$ --- разбиение отрезка $[a, b]$. Выборкой, соответствующей разбиению $T$, будем называть набор точек $\{\xi_i\}_{i = 1}^{N_T}$ такой, что 
    \[\xi_i \in [x_{i - 1}, x_i] \quad \forall i \in \{1, \ldots, N_T\}\]
    Выборка, соответствующая разбиению $T$, обозначается $\xi_T$.
\end{definition}

\begin{definition}
    Пусть $f \colon [a, b] \to \mathbb{R}$. Пусть $T \in \mathcal{T}([a, b])$, а $\xi_T$ --- какая-то выборка, соответствующая разбиению $T$. Тогда интегральной суммой Римана, построенной по разбиению $T$ и выборке $\xi_T$, будем называть следующее:
    \[\sigma(f, T, \xi_T) := \sum_{i = 1}^{N_T} f(\xi_i) (x_i - x_{i - 1})\]
\end{definition}

\begin{theorem}[Критерий интегрируемости в терминах интегральных сумм Римана]
    Пусть  \\
    $f \colon [a, b] \to \mathbb{R}$.
    Следующие условия эквивалентны:
    \begin{enumerate}
        \item $f \in \mathcal{R}([a, b]) \;$ и $\; \displaystyle J = \int_{a}^{b} f(x) \,dx$
        \item $\displaystyle \lim_{l(T) \to 0} \sigma(f, T, \xi_T) = J \iff \forall \varepsilon > 0 \; \exists \delta (\varepsilon) > 0: \forall T \in \mathcal{T}([a, b]): l(T) < \delta(\varepsilon) \text{ и } \forall \text{ выборки } \xi_T \hookrightarrow \left\lvert \sigma(f, T, \xi_T) - J \right\rvert < \varepsilon$
    \end{enumerate}
\end{theorem}

\begin{proof}
    Кванторное условие $2$ эквивалентно следующему:
    \[\forall \varepsilon > 0 \; \exists \delta (\varepsilon) > 0: \forall T \in \mathcal{T}([a, b]): l(T) < \delta(\varepsilon) \text{ и } \forall \xi_T \hookrightarrow  J - \varepsilon \leq \sigma(f, T, \xi_T) \leq J + \varepsilon\]
    \[\iff\]
    \begin{equation*} \label{eq:th:2.1} \tag{$\star$}
        \forall \varepsilon > 0 \; \exists \delta (\varepsilon) > 0: \forall T \in \mathcal{T}([a, b]): l(T) < \delta(\varepsilon) \text{ и } \forall \xi_T \hookrightarrow  J - \varepsilon \leq \sum_{i = 1}^{N_T} f(\xi_i) (x_i - x_{i - 1}) \leq J + \varepsilon
    \end{equation*}
    Рассмотрим неравенства из условия (\ref{eq:th:2.1}) при фиксированных $\varepsilon > 0$ и разбиении $T \in \mathcal{T}([a, b])$ таком, что $l(T) < \delta(\varepsilon)$. Тогда от обеих частей этих неравенств можно взять инфимум и супремум по всем выборкам $\xi_T$. Сделав это, получим следующее:
    \[J - \varepsilon \leq \inf_{\xi_T} \sigma(f, T, \xi_T) \leq J + \varepsilon\] 
    \[J - \varepsilon \leq \sup_{\xi_T} \sigma(f, T, \xi_T) \leq J + \varepsilon\]
    Покажем теперь, что
    \[\inf_{\xi_T} \sigma(f, T, \xi_T) = s_f (T)\]
    Распишем чудесный инфимум в левой части неравенства:
    \[\inf_{\xi_T} \sigma(f, T, \xi_T) = \inf_{\xi_T} \sum_{i = 1}^{N_T} f(\xi_i) (x_i - x_{i - 1}) = \inf_{\substack{\xi_1 \in [x_0, x_1] \\ \vdots \\ \xi_{N_T} \in \left[x_{N_T - 1}, x_{N_T}\right]}} \sum_{i = 1}^{N_T} f(\xi_i) (x_i - x_{i - 1}) = \]
    \[= [\text{лёгкое упражнение}] = \left( \inf_{\xi_1 \in [x_0, x_1]} \ldots \left( \inf_{\xi_{N_T} \in \left[x_{N_T - 1}, x_{N_T}\right]} \sum_{i = 1}^{N_T} f(\xi_i) (x_i - x_{i - 1}) \right) \ldots \right) = \]
    \[= \sum_{i = 1}^{N_T} \left(\inf_{x_{i - 1} \leq \xi \leq x_i} f(\xi)\right) (x_i - x_{i - 1}) = s_f(T)\]
    Аналогично доказывается, что 
    \[\sup_{\xi_T} \sigma(f, T, \xi_T) = S_f (T)\]
    Тогда в итоге при фиксированных $\varepsilon > 0$ и разбиении $T \in \mathcal{T}([a, b])$ таком, что $l(T) < \delta(\varepsilon)$, неравенства (\ref{eq:th:2.1}) просто равносильны следующему:
    \[J - \varepsilon \leq s_f(T) \leq S_f(T) \leq J + \varepsilon\]
    Но $\varepsilon > 0$ и разбиение $T \in \mathcal{T}([a, b])$ такое, что $l(T) < \delta(\varepsilon)$, были выбраны произвольно, откуда:
    \[\forall \varepsilon > 0 \; \exists \delta (\varepsilon) > 0: \forall T \in \mathcal{T}([a, b]): l(T) < \delta(\varepsilon) \text{ и } \forall \xi_T \hookrightarrow J - \varepsilon \leq s_f(T) \leq S_f(T) \leq J + \varepsilon\]
    \[\iff\]
    \[f \in \mathcal{R}([a, b]) \quad \text{ и } \quad J = \int_{a}^{b} f(x) \,dx\]
\end{proof}

\begin{theorem}[Линейность интеграла Римана]
    Множество $\mathcal{R}([a, b])$ является линейным пространством, то есть
    \[\forall \alpha, \beta \in \mathbb{R}, \; \forall f, g \in \mathcal{R}([a, b]) \hookrightarrow \alpha f + \beta g \in \mathcal{R}([a, b])\]
    Более того, 
    \[\int_{a}^{b} (\alpha f + \beta g)(x) \,dx = \alpha \int_{a}^{b} f(x) \,dx + \beta \int_{a}^{b} g(x) \,dx\]
\end{theorem}

\begin{proof}
    Пусть $T \in \mathcal{T}([a, b])$. Пусть $\xi_T$ --- произвольная выборка, соответствующая разбиению $T$. Очевидно, что
    \begin{equation*} \label{eq:th:2.2} \tag{$\star \star$}
        \sigma(\alpha f + \beta g, T, \xi_T) = \alpha \sigma(f, T, \xi_T) + \beta \sigma (g, T, \xi_T)
    \end{equation*}
    Будем считать, что $\left\lvert \alpha \right\rvert + \left\lvert \beta \right\rvert > 0$, так как иначе всё тривиально. Так как $f, g \in \mathcal{R}([a, b])$, то по критерию интегрируемости в терминах интегральных сумм Римана
    \[\exists J_1 = \int_{a}^{b} f(x) \,dx \in \mathbb{R}, \; \; \exists J_2 =  \int_{a}^{b} g(x) \,dx \in \mathbb{R} : \]
    \[\forall \varepsilon > 0 \; \exists \delta_1(\varepsilon) > 0: \forall T \in \mathcal{T}([a, b]) : l(T) < \delta_1(\varepsilon) \text{ и } \forall \xi_T \hookrightarrow \left\lvert J_1 - \sigma(f, T, \xi_T) \right\rvert < \frac{\varepsilon}{2 \left( \left\lvert \alpha \right\rvert + \left\lvert \beta \right\rvert \right)}\]
    \[\forall \varepsilon > 0 \; \exists \delta_2(\varepsilon) > 0: \forall T \in \mathcal{T}([a, b]) : l(T) < \delta_2(\varepsilon) \text{ и } \forall \xi_T \hookrightarrow \left\lvert J_2 - \sigma(g, T, \xi_T) \right\rvert < \frac{\varepsilon}{2 \left( \left\lvert \alpha \right\rvert + \left\lvert \beta \right\rvert \right)}\]
    Положим $\delta(\varepsilon) := \min \{\delta_1(\varepsilon), \delta_2(\varepsilon)\}$ и, суммируя сказанное выше, получим:
    \[\forall \varepsilon > 0 \; \exists \delta(\varepsilon) > 0: \forall T \in \mathcal{T}([a, b]): l(T) < \delta(\varepsilon) \text{ и } \forall \xi_T \hookrightarrow \left\lvert \sigma(\alpha f + \beta g, T, \xi_T) - \alpha J_1 - \beta J_2 \right\rvert \leq \]
    \[\leq [\text{неравенство треугольника + (\ref{eq:th:2.2})}] \leq \left\lvert \alpha \right\rvert \cdot \left\lvert J_1 - \sigma (f, T, \xi_T) \right\rvert + \left\lvert \beta \right\rvert \cdot \left\lvert J_2 - \sigma(g, T, \xi_T) \right\rvert < \]
    \[< \frac{\varepsilon \left\lvert \alpha \right\rvert}{2 \left( \left\lvert \alpha \right\rvert + \left\lvert \beta \right\rvert \right)} + \frac{\varepsilon \left\lvert \beta \right\rvert}{2 \left( \left\lvert \alpha \right\rvert + \left\lvert \beta \right\rvert \right)} \leq \frac{\varepsilon}{2} + \frac{\varepsilon}{2} = \varepsilon\]
    Из этого по критерию интегрируемости в терминах интегральных сумм Римана немедленно получаем, что 
    \[\alpha f + \beta g \in \mathcal{R}([a, b]) \quad \text{ и } \quad \int_{a}^{b} (\alpha f + \beta g)(x) \,dx = \alpha \int_{a}^{b} f(x) \,dx + \beta \int_{a}^{b} g(x) \,dx\]
    Легко видеть, что требуемое доказано.
\end{proof}

\end{document}